%!TEX root = ../thesis.tex
\chapter{Experiments and Results}
\label{chap:experiments}

%% This chapter presents all experimental results.
%% Structure: 4 phases matching the methodology, plus "Why SSL Works" analysis.
%% Approx. 20 pages — the core of the thesis.

% ============================================================================
\section{Baseline Validation}
\label{sec:exp-baselines}

Before interpreting cross-city behavior, we first establish the scale of the task on the standard mixed-city NAVSIM split. Table~\ref{tab:baselines} includes both non-learning or low-capacity lower bounds and strong learned baselines. The purpose is not to reproduce the full leaderboard, but to anchor the later discussion of representation quality against concrete reference points.

\begin{table}[ht]
\centering
\caption{Reference baselines on NAVSIM under the standard mixed-city split.}
\label{tab:baselines}
\begin{tabular}{llc}
\toprule
\textbf{Exp} & \textbf{Configuration} & \textbf{PDMS} \\
\midrule
A1  & Constant-velocity controller & 25.88 \\
A2  & Ego-status MLP & 65.99 \\
A4  & TransFuser (ResNet34, C+L) & 81.68 \\
A10 & DiffusionDrive (ResNet34, C+L) & 88.1 \\
\bottomrule
\end{tabular}
\end{table}

These references establish two points used throughout the rest of the chapter. First, the benchmark is non-trivial: ego-state prediction alone is insufficient. Second, planner architecture can clearly improve mixed-city PDMS, as seen in the jump from TransFuser to DiffusionDrive. The thesis question is therefore more specific than which planner scores highest in aggregate. It is which parts of that performance survive under geographic transfer.

% ============================================================================
\section{Exploratory Lightweight-Head Study}
\label{sec:exp-phase0}

This thesis did not begin with city splits. The original working hypothesis was simpler: if representation quality is genuinely central to end-to-end driving, then replacing a supervised image backbone with a self-supervised one should already help even when the planner head is intentionally weak. The exploratory regime of~\ref{sec:meth-lightweight-head} operationalizes this hypothesis through the \texttt{LightweightPlanningAgent}: eleven backbones, two adaptation modes, the same MLP head, the same training recipe, the same 20-epoch budget on the pooled \texttt{navtest} split. Its purpose is not leaderboard ambition but causal cleanliness.

\subsection{Priority-1 Backbone Sweep (Frozen, Front-1, MLP)}
\label{sec:exp-phase0-p1}

The Priority-1 sweep holds the planner at its minimum capacity --- a single front camera, frozen backbone, $1288{\to}256{\to}256{\to}T{\cdot}3$ MLP --- and varies only the backbone. \reftab{tab:exp-phase0-p1} reports all ten Priority-1 results on the full \texttt{navtest} split, ranked by PDMS. Each row is one evaluation record in the authoritative MongoDB experiment store, and sub-scores (NC, DAC, EP, TTC, C) come from the same record.

\begin{table}[ht]
\centering
\caption{Priority-1 lightweight-head sweep on \texttt{navtest} ($12{,}146$ scenarios), pooled all-city. All rows use the same frozen encoder, single-front-camera, MLP$(256,256)$ head, AdamW, 20-epoch recipe from~\ref{sec:meth-lh-recipe}. PDMS is one-stage, devkit v1.0, non-reactive.}
\label{tab:exp-phase0-p1}
\small
\renewcommand{\arraystretch}{1.15}
\setlength{\tabcolsep}{4pt}
\begin{tabular}{@{}>{\raggedright\arraybackslash}p{3.2cm} >{\raggedright\arraybackslash}p{2.4cm} >{\raggedright\arraybackslash}p{2.6cm} ccccccc@{}}
\toprule
\textbf{Experiment} & \textbf{Backbone} & \textbf{Pre-training} & \textbf{PDMS} & \textbf{NC} & \textbf{DAC} & \textbf{EP} & \textbf{TTC} & \textbf{C} \\
\midrule
\texttt{LH-B1-IJH-hf-fr.}  & I-JEPA ViT-H/14 & ImageNet-1k SSL         & $\mathbf{0.797}$ & 0.966 & 0.892 & 0.868 & 0.954 & 0.935 \\
\texttt{LH-B1-MAE-hf-fr.}  & MAE ViT-B/16    & ImageNet-1k SSL         & $0.775$ & 0.962 & 0.871 & 0.869 & 0.949 & 0.937 \\
\texttt{LH-B1-MAE-s-fr.}   & MAE ViT-S/14    & nuSc $224^2$            & $0.749$ & 0.948 & 0.856 & 0.870 & 0.935 & 0.937 \\
\texttt{LH-B1-MAE-r-fr.}   & MAE ViT-S/14    & nuSc $224{\times}392$   & $0.738$ & 0.948 & 0.846 & 0.868 & 0.935 & 0.934 \\
\texttt{LH-B1-IJS-s-fr.}   & I-JEPA ViT-S/14 & nuSc $224^2$            & $0.716$ & 0.941 & 0.827 & 0.868 & 0.928 & 0.936 \\
\texttt{LH-B1-Dv2-r-fr.}   & DINOv2 ViT-S/14 & nuSc $224{\times}392$   & $0.698$ & 0.942 & 0.809 & 0.861 & 0.927 & 0.929 \\
\texttt{LH-B1-Dv2-s-fr.}   & DINOv2 ViT-S/14 & nuSc $224^2$            & $0.697$ & 0.942 & 0.807 & 0.861 & 0.927 & 0.933 \\
\rowcolor{gray!15}
\texttt{LH-B1-Rn34-fr.}    & ResNet-34       & ImageNet supervised     & $0.687$ & 0.936 & 0.798 & 0.866 & 0.921 & 0.927 \\
\texttt{LH-B1-Dv2-hf-fr.}  & DINOv2 ViT-S/14 & LVD-142M SSL            & $0.668$ & 0.934 & 0.781 & 0.858 & 0.918 & 0.931 \\
\texttt{LH-B1-IJS-r-fr.}   & I-JEPA ViT-S/14 & nuSc $224{\times}392$   & $0.654$ & 0.931 & 0.767 & 0.857 & 0.914 & 0.928 \\
\bottomrule
\end{tabular}
\end{table}

\subsection{Priority-3 Sweep (Fully Fine-Tuned, Front-1, MLP)}
\label{sec:exp-phase0-p3}

Priority-3 repeats the same backbone matrix with \(\texttt{trainable\_backbone\_fraction}{=}1.0\) and the two-group AdamW schedule (head LR $10^{-4}$, encoder LR $3 \times 10^{-5}$) from~\ref{sec:meth-lh-recipe}. Nine of ten rows are scored; the I-JEPA ViT-H/14 fully-unfrozen run is re-evaluating after an out-of-memory event on the one-stage scorer and is marked pending.

\begin{table}[ht]
\centering
\caption{Priority-3 lightweight-head sweep --- fully unfrozen backbone, otherwise identical to Table~\ref{tab:exp-phase0-p1}.}
\label{tab:exp-phase0-p3}
\small
\renewcommand{\arraystretch}{1.15}
\setlength{\tabcolsep}{4pt}
\begin{tabular}{@{}>{\raggedright\arraybackslash}p{3.2cm} >{\raggedright\arraybackslash}p{2.4cm} >{\raggedright\arraybackslash}p{2.6cm} ccccccc@{}}
\toprule
\textbf{Experiment} & \textbf{Backbone} & \textbf{Pre-training} & \textbf{PDMS} & \textbf{NC} & \textbf{DAC} & \textbf{EP} & \textbf{TTC} & \textbf{C} \\
\midrule
\texttt{LH-B1-MAE-hf-full} & MAE ViT-B/16    & ImageNet-1k SSL       & $\mathbf{0.849}$ & 0.974 & 0.936 & 0.869 & 0.963 & 0.939 \\
\texttt{LH-B1-MAE-r-full}  & MAE ViT-S/14    & nuSc $224{\times}392$ & $0.817$ & 0.966 & 0.914 & 0.868 & 0.955 & 0.936 \\
\texttt{LH-B1-MAE-s-full}  & MAE ViT-S/14    & nuSc $224^2$          & $0.814$ & 0.961 & 0.914 & 0.870 & 0.951 & 0.937 \\
\texttt{LH-B1-IJS-s-full}  & I-JEPA ViT-S/14 & nuSc $224^2$          & $0.809$ & 0.964 & 0.906 & 0.871 & 0.952 & 0.939 \\
\rowcolor{gray!15}
\texttt{LH-B1-Rn34-full}   & ResNet-34       & ImageNet supervised   & $0.769$ & 0.959 & 0.866 & 0.871 & 0.947 & 0.935 \\
\texttt{LH-B1-Dv2-s-full}  & DINOv2 ViT-S/14 & nuSc $224^2$          & $0.760$ & 0.954 & 0.865 & 0.868 & 0.941 & 0.929 \\
\texttt{LH-B1-Dv2-r-full}  & DINOv2 ViT-S/14 & nuSc $224{\times}392$ & $0.760$ & 0.954 & 0.865 & 0.868 & 0.941 & 0.929 \\
\texttt{LH-B1-Dv2-hf-full} & DINOv2 ViT-S/14 & LVD-142M SSL          & $0.718$ & 0.943 & 0.826 & 0.866 & 0.929 & 0.930 \\
\texttt{LH-B1-IJS-r-full}  & I-JEPA ViT-S/14 & nuSc $224{\times}392$ & $0.668$ & 0.932 & 0.780 & 0.861 & 0.918 & 0.929 \\
\texttt{LH-B1-IJH-hf-full} & I-JEPA ViT-H/14 & ImageNet-1k SSL       & \emph{pending} & -- & -- & -- & -- & -- \\
\bottomrule
\end{tabular}
\end{table}

\subsection{Adaptation Lift: Frozen $\to$ Fully Unfrozen}
\label{sec:exp-phase0-lift}

Pairing each row of Table~\ref{tab:exp-phase0-p1} with its Table~\ref{tab:exp-phase0-p3} counterpart gives the adaptation lift $\Delta_\text{adapt}$ per backbone.

\begin{table}[ht]
\centering
\caption{Per-backbone lift from fully unfreezing the encoder at Priority-3 relative to the frozen Priority-1 baseline.}
\label{tab:exp-phase0-lift}
\small
\renewcommand{\arraystretch}{1.15}
\setlength{\tabcolsep}{6pt}
\begin{tabular}{@{}>{\raggedright\arraybackslash}p{4.8cm} c c c@{}}
\toprule
\textbf{Backbone variant} & \textbf{Frozen} & \textbf{Full} & $\boldsymbol{\Delta_\text{adapt}}$ \\
\midrule
MAE ViT-S/14 (nuSc $224^2$)        & $0.749$ & $0.814$ & $+0.064$ \\
MAE ViT-S/14 (nuSc $224{\times}392$) & $0.738$ & $0.817$ & $+0.078$ \\
MAE ViT-B/16 (ImageNet)            & $0.775$ & $0.849$ & $+0.074$ \\
I-JEPA ViT-S/14 (nuSc $224^2$)     & $0.716$ & $0.809$ & $\mathbf{+0.094}$ \\
I-JEPA ViT-S/14 (nuSc $224{\times}392$) & $0.654$ & $0.668$ & $+0.014$ \\
DINOv2 ViT-S/14 (nuSc $224^2$)     & $0.697$ & $0.760$ & $+0.063$ \\
DINOv2 ViT-S/14 (nuSc $224{\times}392$) & $0.698$ & $0.760$ & $+0.062$ \\
DINOv2 ViT-S/14 (LVD-142M)         & $0.668$ & $0.718$ & $+0.050$ \\
\rowcolor{gray!15}
ResNet-34 (ImageNet, supervised)   & $0.687$ & $0.769$ & $+0.082$ \\
\bottomrule
\end{tabular}
\end{table}

\subsection{Cross-City Probe (Priority-2 CC)}
\label{sec:exp-phase0-cc}

The cross-city leg of the Exploratory matrix trains four frozen backbones on each of the four individual NAVSIM cities and evaluates every resulting checkpoint on the full \texttt{navtest} split. All four backbones use their ImageNet-pretrained release (the \texttt{-hf} variants) through the custom ViT backend, so capacity is not matched across rows but pre-training domain is held constant. This is the minimum-capacity analogue of the TransFuser / Latent TransFuser cross-city tables in our earlier study~\cite{naeinian-cross-city-2026}, and it is carried forward into the cross-city discussion in~\ref{sec:exp-phase1}.

\begin{table}[ht]
\centering
\caption{Exploratory Cross-City probe: PDMS on full \texttt{navtest} for models trained on each single city, front-1 frozen MLP head. The last row is the per-backbone mean across the four training cities.}
\label{tab:exp-phase0-cc}
\small
\renewcommand{\arraystretch}{1.2}
\setlength{\tabcolsep}{6pt}
\begin{tabular}{@{}>{\raggedright\arraybackslash}p{3.0cm} cccc@{}}
\toprule
\textbf{Train city} & \textbf{ResNet-34} & \textbf{DINOv2 ViT-S/14} & \textbf{I-JEPA ViT-H/14} & \textbf{MAE ViT-B/16} \\
\midrule
Boston      & $0.611$ & $0.598$ & $\mathbf{0.688}$ & $0.654$ \\
Pittsburgh  & $0.556$ & $0.557$ & $\mathbf{0.652}$ & $0.613$ \\
Singapore   & $0.511$ & $0.509$ & $\mathbf{0.529}$ & $0.528$ \\
Las Vegas   & $0.623$ & $0.595$ & $\mathbf{0.722}$ & $0.692$ \\
\midrule
\rowcolor{gray!15}
Mean (4 cities) & $0.575$ & $0.564$ & $\mathbf{0.648}$ & $0.622$ \\
\bottomrule
\end{tabular}
\end{table}

\subsection{Findings}
\label{sec:exp-phase0-findings}

Seven observations follow from the Priority-1, Priority-3, and Cross-City results above.

\paragraph{Frozen I-JEPA ViT-H/14 is the strongest minimum-head configuration.} At \(0.797\) PDMS with zero backbone gradients updated, \texttt{LH-B1-IJH-hf-frozen} (Table~\ref{tab:exp-phase0-p1}) outperforms every fully-unfrozen small-ViT configuration in Table~\ref{tab:exp-phase0-p3} except the three MAE runs. This is the strongest single argument for frozen self-supervised features plus a lightweight adapter: the representation carries enough planning-relevant structure that backbone gradients do not have to do the work.

\paragraph{MAE leads the fully-fine-tuned leaderboard; I-JEPA leads the frozen one.} Under Priority-3, the three MAE rows (\(0.849\), \(0.817\), \(0.814\)) hold the top three positions and together exceed every I-JEPA or DINOv2 variant. Under Priority-1, I-JEPA ViT-H/14 leads by \(2.2\) points over the best MAE configuration. The reconstructive objective of MAE appears to produce features that adapt better under end-to-end fine-tuning, while I-JEPA's latent-predictive features are more usable off-the-shelf. Both are consistent with the thesis claim that the SSL objective is the variable that matters.

\paragraph{Self-supervised pretraining beats supervised ResNet-34 under the same head.} In the frozen setting, seven of nine SSL variants exceed the ResNet-34 baseline (\(0.687\)); under full fine-tuning, seven of nine again exceed ResNet-34 (\(0.769\)). The two SSL rows that trail are the off-the-shelf DINOv2 LVD-142M checkpoint and the I-JEPA ViT-S/14 rectangular variant --- both architectural edge cases rather than typical deployments.

\paragraph{Fine-tuning lift depends on the backbone.} Table~\ref{tab:exp-phase0-lift} shows I-JEPA ViT-S/14 (square) is the biggest beneficiary of fine-tuning ($+9.4$\,pp), followed by supervised ResNet-34 ($+8.2$\,pp) and the two MAE ViT-S/14 variants ($+7.8$\,pp and $+6.4$\,pp). The exception is \textbf{I-JEPA ViT-S/14 rectangular}, which lifts only $+1.4$\,pp. This flags a likely optimizer-instability issue when the EMA target-encoder weights of a rectangular pretraining are unfrozen together with their non-square positional embeddings; the thesis treats this row as a known-flaky configuration rather than as evidence against I-JEPA.

\paragraph{Rectangular pretraining does not help pooled in-domain PDMS.} For both MAE and I-JEPA, the square ViT-S/14 variant outperforms its rectangular counterpart on pooled \texttt{navtest} in both adaptation modes. DINOv2 is tied within $0.001$. This replicates the observation in our earlier benchmark study~\cite{naeinian-cross-city-2026} that the rectangular-input advantage is specific to explicit cross-city evaluation and does not appear on pooled in-domain splits, and it motivates the inversion of the comparison in~\ref{sec:exp-phase1}.

\paragraph{Cross-city ranking follows the pooled ranking but amplifies the SSL gap.} Averaged across the four single-city training sources, I-JEPA ViT-H/14 (\(0.648\)) and MAE ViT-B/16 (\(0.622\)) both beat supervised ResNet-34 (\(0.575\)) by \(7.3\) and \(4.7\) points respectively. DINOv2 (\(0.564\)) trails ResNet-34 slightly, consistent with its weaker pooled behavior. The representation advantage becomes larger, not smaller, once the evaluation distribution is a different city from the training distribution --- this is the Exploratory preview of the main thesis claim.

\paragraph{Singapore is the cross-city anomaly.} All four backbones fall to their lowest PDMS when trained on Singapore (\(0.509\)--\(0.529\)), a spread of only \(2\) PDMS points across the backbone families. Singapore has the smallest single-city partition (\(97\) scenes, \(5.8\) hours) and is the only left-hand-traffic city; the representation advantage is partially eaten by the limited gradient signal available during downstream training. This matches the directional-asymmetry finding in our earlier cross-city study~\cite{naeinian-cross-city-2026} exactly.

\subsection{In-Flight B-series Ablations}
\label{sec:exp-phase0-inflight}

The nine B-series runs (multi-camera and fusion/head ablations from~\ref{sec:meth-lh-taxonomy}) are currently queued behind the rebuild of the \texttt{front3} and \texttt{front5\_plus\_b} training caches. \reftab{tab:exp-phase0-inflight} lists the block for completeness; all rows share the I-JEPA ViT-H/14 ImageNet backbone and are scheduled to fire automatically once the squashfs compressions complete.

\begin{table}[ht]
\centering
\caption{Exploratory B-series ablations (queued, I-JEPA ViT-H/14 backbone throughout).}
\label{tab:exp-phase0-inflight}
\small
\renewcommand{\arraystretch}{1.15}
\setlength{\tabcolsep}{5pt}
\begin{tabular}{@{}>{\raggedright\arraybackslash}p{2.9cm} >{\raggedright\arraybackslash}p{2.2cm} >{\raggedright\arraybackslash}p{2.6cm} >{\raggedright\arraybackslash}p{1.9cm} >{\raggedright\arraybackslash}p{2.0cm}@{}}
\toprule
\textbf{Experiment} & \textbf{Cameras} & \textbf{Fusion} & \textbf{Head} & \textbf{Trainable} \\
\midrule
\texttt{LH-B16}        & front3      & mean            & MLP         & frozen \\
\texttt{LH-B17-frozen} & front5+b    & mean            & MLP         & frozen \\
\texttt{LH-B17-full}   & front5+b    & mean            & MLP         & full \\
\texttt{LH-B19-frozen} & front3      & cross\_attn     & MLP         & frozen \\
\texttt{LH-B19-full}   & front3      & cross\_attn     & MLP         & full \\
\texttt{LH-B20-frozen} & front5+b    & cross\_attn     & MLP         & frozen \\
\texttt{LH-B20-full}   & front5+b    & cross\_attn     & MLP         & full \\
\texttt{LH-B21}        & front3      & weighted\_mean  & MLP         & full \\
\texttt{LH-B23}        & front5+b    & weighted\_mean  & Transformer & full \\
\bottomrule
\end{tabular}
\end{table}

\subsection{Drive-JEPA and the Post-Saturation Reframe}
\label{sec:exp-phase0-drivejepa}

The spread in Tables~\ref{tab:exp-phase0-p1}--\ref{tab:exp-phase0-p3} is consistent with the thesis claim that representations dominate planner capacity in the lightweight regime, but it is not by itself a novel contribution on pooled NAVSIM. Drive-JEPA~\cite{wang-drive-jepa} reports that a V-JEPA video encoder paired with a simple transformer decoder reaches $89.0$ PDMS on the same benchmark, above every single-image SSL configuration in these two tables (top-row: frozen I-JEPA ViT-H/14 at $0.797$, fully-fine-tuned MAE ViT-B/16 at $0.849$), and the full Drive-JEPA system reaches $93.3$ PDMS on NAVSIM v1 and $87.8$ EPDMS on v2. That result validates the representation-centric direction behind the exploratory study but makes raw pooled PDMS an insufficient differentiator for the thesis as a whole.

The exploratory study therefore plays the role stated in~\ref{sec:meth-lh-role}: a diagnostic substrate that establishes what a minimum-capacity planner can discriminate about the backbone, with the city-disjoint evaluation in~\ref{sec:exp-phase1} serving as the main vehicle for the novel claim. The remainder of the chapter is structured around that pivot: explicit cross-city evaluation (the cross-city benchmark, \refsec{sec:exp-phase1}), architectural transfer of the representation effect to a generative planner (the generative-planning studies, \refsec{sec:exp-phase2}--\ref{sec:exp-phase3}), and a mechanistic account of why SSL features transfer more gracefully (the mechanistic analysis, \refsec{sec:exp-why-ssl}).

% ============================================================================
\section{Cross-City SSL Benchmark}
\label{sec:exp-phase1}

This section synthesizes the core cross-city benchmark reported in our earlier study~\cite{naeinian-cross-city-2026}. It establishes the empirical foundation for the thesis and explains the pivot away from pure mixed-city performance. Once geographic mixing is removed from evaluation, robustness depends much more strongly on the visual representation than on the downstream planning head.

\subsection{Experimental Setup}
The phase-1 protocol has two complementary parts. In open-loop evaluation, LAW is trained on nuScenes under strict Boston/Singapore geographic splits and evaluated in both in-domain and cross-city directions. In closed-loop evaluation, TransFuser and Latent TransFuser are trained on single cities in NAVSIM and evaluated on the remaining three cities. Across both settings, the downstream planner is held fixed while the backbone varies across supervised baselines, generic ImageNet-pretrained SSL models, and domain-specific nuScenes-pretrained I-JEPA, DINOv2, and MAE variants. Frozen and fully fine-tuned downstream settings are both included, which makes it possible to separate the value of the pretrained representation from the value of downstream adaptation~\cite{naeinian-cross-city-2026}.

\subsection{In-Distribution Performance}
\label{sec:exp-phase1-indist}

The in-distribution story is intentionally unspectacular. On the standard mixed-city nuScenes split, the supervised Swin baseline reaches 0.69 average L2 and 0.22 average collision rate, while strong SSL alternatives remain in the same range: I-JEPA ViT-H/14 fine-tuned reaches 0.63 / 0.28, and the nuScenes-pretrained rectangular I-JEPA ViT-S/14 frozen model reaches 0.61 / 0.35~\cite{naeinian-cross-city-2026}. This is the pattern that motivates the rest of the chapter: mixed-city validation suggests that most competent backbones are roughly comparable, so it hides the degree to which some of them are learning city-specific shortcuts.

The same phenomenon appears in closed-loop evaluation. Under all-city NAVSIM training, performance differences across backbones are relatively small because the optimizer sees geographically diverse supervision. In other words, standard validation mostly measures interpolation inside the support of the pooled dataset, not robustness to a true city shift.

\subsection{Cross-City Transfer Matrices}
\label{sec:exp-phase1-matrix}

The transfer matrices reveal structure that is invisible in the pooled metrics. On nuScenes, the Boston$\rightarrow$Singapore direction is consistently harder than Singapore$\rightarrow$Boston, especially for supervised or weakly adapted backbones. On NAVSIM, the full four-city matrix shows the same phenomenon at larger scale: Boston is generally a stronger training city than Singapore for single-city transfer, while Singapore is the harshest test environment because it combines the smallest dataset with the only left-hand-traffic regime~\cite{naeinian-cross-city-2026}. The resulting pattern answers RQ2 in a practical way: the ``best'' source city is the one that exposes the model to structural diversity without overfitting it to idiosyncratic local cues.

\begin{figure}[ht]
    \centering
    \includegraphics[width=0.48\textwidth]{iros/fig1_tf_pdms_heatmaps.pdf}
    \includegraphics[width=0.48\textwidth]{iros/fig1_latf_pdms_heatmaps.pdf}
    \caption{Four-city cross-city PDMS matrices for TransFuser (left) and Latent TransFuser (right), pulled directly from our earlier cross-city benchmark study~\cite{naeinian-cross-city-2026}. Rows are training cities, columns are test cities, and each panel is a backbone family. The diagonal is in-distribution; off-diagonal cells expose the representation-dependent generalization gap. For every backbone, the Singapore column is the harshest; for supervised ResNet-34 (top-left panel), the off-diagonal penalty is larger and more uneven than for the frozen SSL panels.}
    \label{fig:iros-pdms-heatmaps}
\end{figure}

\subsection{OOD PDMS Analysis}
\label{sec:exp-phase1-ood}

The most direct open-loop evidence comes from transfer from Boston to Singapore on nuScenes. The supervised Swin backbone inflates from 0.60 to 5.86 average L2 and from 0.32 to 6.19 average collision rate, corresponding to 9.77$\times$ and 19.34$\times$ degradation. Replacing the backbone with a domain-specific SSL representation changes the picture entirely: the rectangular, frozen nuScenes-pretrained I-JEPA ViT-S/14 model yields an error ratio of only 1.20$\times$ in L2 and 0.75$\times$ in collision rate for the same transfer~\cite{naeinian-cross-city-2026}. The planning architecture is unchanged, so the improved robustness must come from the representation.

The closed-loop NAVSIM results point in the same direction. Generic large-scale SSL often matches or slightly exceeds the supervised ResNet34 baseline under single-city transfer, and domain-specific SSL is more reliable still. In the LiDAR-free Latent TransFuser setting, nuScenes-pretrained MAE improves average OOD PDMS by more than four points for Boston-, Pittsburgh-, and Las Vegas-trained models and by roughly three points for Singapore-trained models~\cite{naeinian-cross-city-2026}. The effect is therefore not a quirk of one backbone family or one planner: once city shift is explicit, self-supervised initialization systematically narrows the gap between in-domain and out-of-domain behavior.

\begin{figure}[ht]
    \centering
    \includegraphics[width=0.75\textwidth]{iros/fig3_indomain_vs_ood.pdf}
    \caption{In-distribution PDMS versus average out-of-distribution PDMS across every cross-city configuration. Each point is one trained model; the diagonal is the reference ``no transfer penalty'' line. Points above the diagonal indicate in-distribution overfitting relative to OOD behavior. Supervised backbones (red) cluster further from the diagonal than frozen SSL backbones (blue) at every in-distribution score, consistent with the representation hypothesis. Source: Fig.~3 of our earlier cross-city benchmark study~\cite{naeinian-cross-city-2026}.}
    \label{fig:iros-indomain-vs-ood}
\end{figure}

Closed-loop cross-city transfer is further evaluated using the NAVSIM benchmark with the Performance Degradation Metric Score (PDMS), which measures driving task completion while penalizing infractions (higher is better). Table~\ref{tab:navsim_pdms_merged} reports PDMS across four cities --- Las Vegas (LV), Boston (B), Pittsburgh (P), and Singapore (S) --- for both TransFuser and Latent TransFuser under the NAVSIM v2.2 devkit; Table~\ref{tab:navsim_pdms_merged_v11} is the same matrix re-evaluated with the official NAVSIM v1.1 devkit as a defensive check against devkit-version confounders. For each backbone, we report results from all-city and single-city training to illustrate the impact of the geographic domain shift on closed-loop performance.

% Full NAVSIM v2.2 cross-city PDMS matrix (legacy devkit).
% Retained as provenance for the companion paper's numbers; the canonical
% thesis matrix is the v1.1 table (tab_navsim_pdms_merged_v11.tex).
\begin{sidewaystable}[p]
    \centering
    \caption[NAVSIM v2.2 cross-city PDMS matrix (legacy).]{%
    Full closed-loop cross-city PDMS matrix on NAVSIM, devkit v2.2 (legacy).
    Each row is a single trained model. The left block (\emph{Train City}) marks which of the four cities were used for training; \cmark indicates inclusion. Total combines performance across all four evaluation cities for each planner. All numbers are PDMS (\%); higher is better.
    How to read it: rows with four training cities are the all-city baselines, and their \emph{Total} column is mixed-city interpolation performance. Rows with a single training city are zero-shot transfer probes; non-checked columns under each planner are the out-of-distribution cells, and the gap between the in-distribution column and the three OOD columns quantifies the geographic transfer penalty. This matrix is retained for direct comparison with the companion cross-city benchmark paper~\cite{naeinian-cross-city-2026}; the thesis claims use the v1.1 re-evaluation in~\reftab{tab:navsim_pdms_merged_v11}.}
    \label{tab:navsim_pdms_merged}
    \scriptsize
    \setlength{\tabcolsep}{2pt}
    \renewcommand{\arraystretch}{0.95}
    \resizebox{\textwidth}{!}{%
    \begin{tabular}{|l|cccc|ccccc|ccccc|}
    \toprule
    \multirow{2}{*}{\textbf{Backbone}}
    & \multicolumn{4}{c|}{\textbf{Train City}}
    & \multicolumn{5}{c|}{\textbf{TransFuser (PDMS \%)}}
    & \multicolumn{5}{c|}{\textbf{Latent TransFuser (PDMS \%)}} \\
    \cmidrule(lr){2-5}\cmidrule(lr){6-10}\cmidrule(lr){11-15}
    & LV & B & P & S
    & LV & B & P & S & Total
    & LV & B & P & S & Total \\
    \midrule
    ResNet-34 supervised          & \cmark & \cmark & \cmark & \cmark & 88.4 & \textbf{81.4} & \textbf{75.7} & \textbf{59.9} & \textbf{79.2} & 87.0 & \textbf{79.3} & \textbf{74.4} & \textbf{67.2} & \textbf{79.0} \\
    ResNet-34 supervised          & --     & \cmark & --     & --     & 75.3 & 73.3 & 53.6 & 41.5 & 65.0 & 67.9 & 72.2 & 50.0 & 44.0 & 61.8 \\
    ResNet-34 supervised          & \cmark & --     & --     & --     & \textbf{89.2} & 67.0 & 55.1 & 45.3 & 68.6 & \textbf{87.2} & 63.6 & 50.9 & 41.9 & 65.5 \\
    ResNet-34 supervised          & --     & --     & \cmark & --     & 68.8 & 55.3 & 52.9 & 43.0 & 57.4 & 68.2 & 58.0 & 51.7 & 44.3 & 58.0 \\
    ResNet-34 supervised          & --     & --     & --     & \cmark & 57.9 & 45.5 & 43.4 & 41.0 & 48.5 & 59.3 & 43.9 & 41.1 & 44.6 & 48.6 \\
    \midrule
    I-JEPA ViT-H/14 frozen        & \cmark & \cmark & \cmark & \cmark & 86.5 & 77.1 & 71.2 & 68.1 & 77.6 & 85.0 & 76.5 & 69.6 & 68.7 & 76.7 \\
    I-JEPA ViT-H/14 100\% trainable & \cmark & \cmark & \cmark & \cmark & 87.4 & \textbf{77.7} & 68.8 & 64.2 & 77.0 & \textbf{89.3} & \textbf{77.8} & \textbf{73.0} & 67.6 & \textbf{79.1} \\
    I-JEPA ViT-S/14 nuSc rect., frozen        & \cmark & \cmark & \cmark & \cmark & 88.4 & 77.5 & \textbf{71.5} & 68.4 & \textbf{78.5} & 87.6 & 76.5 & 71.1 & 57.0 & 76.1 \\
    I-JEPA ViT-S/14 nuSc rect., trainable     & \cmark & \cmark & \cmark & \cmark & 88.9 & \textbf{77.7} & 70.8 & 67.0 & 78.4 & 86.6 & 74.8 & 67.2 & 65.3 & 75.7 \\
    I-JEPA ViT-S/14 nuSc sq., frozen          & \cmark & \cmark & \cmark & \cmark & 85.5 & 75.9 & 68.2 & \textbf{69.4} & 76.6 & 84.5 & 73.3 & 68.0 & 65.1 & 74.7 \\
    I-JEPA ViT-S/14 nuSc sq., trainable       & \cmark & \cmark & \cmark & \cmark & 86.7 & 76.2 & 70.9 & 65.7 & 77.0 & 82.9 & 69.8 & 59.1 & \textbf{70.2} & 72.1 \\
    I-JEPA ViT-H/14 frozen        & --     & \cmark & --     & --     & 74.1 & 71.2 & 54.2 & 48.8 & 65.2 & 70.2 & 66.4 & 54.7 & 41.2 & 61.3 \\
    I-JEPA ViT-H/14 trainable     & --     & \cmark & --     & --     & 72.5 & 68.0 & 52.4 & 40.9 & 62.1 & 71.7 & 68.5 & 55.0 & 47.8 & 63.6 \\
    I-JEPA ViT-H/14 frozen        & \cmark & --     & --     & --     & 88.4 & 63.6 & 52.5 & 50.5 & 67.6 & 87.2 & 59.2 & 50.5 & 46.4 & 64.8 \\
    I-JEPA ViT-H/14 trainable     & \cmark & --     & --     & --     & 88.3 & 62.2 & 52.9 & 50.1 & 67.1 & 83.1 & 62.0 & 52.5 & 44.7 & 64.4 \\
    I-JEPA ViT-H/14 frozen        & --     & --     & \cmark & --     & 72.1 & 62.2 & 62.5 & 47.4 & 63.3 & 72.0 & 62.9 & 63.9 & 47.5 & 63.7 \\
    I-JEPA ViT-H/14 trainable     & --     & --     & \cmark & --     & 70.1 & 60.4 & 60.1 & 37.4 & 60.0 & 64.0 & 60.1 & 55.8 & 33.8 & 56.4 \\
    I-JEPA ViT-H/14 frozen        & --     & --     & --     & \cmark & 35.1 & 30.5 & 34.1 & 33.4 & 33.2 & 59.6 & 45.1 & 38.8 & 53.9 & 50.1 \\
    I-JEPA ViT-H/14 trainable     & --     & --     & --     & \cmark & 58.1 & 44.9 & 43.8 & 48.8 & 49.7 & 64.0 & 44.9 & 41.3 & 50.8 & 51.5 \\
    \midrule
    DINOv2 ViT-S/14 frozen                    & \cmark & \cmark & \cmark & \cmark & 88.7 & 78.4 & 70.2 & 66.1 & 78.3 & 87.6 & 76.5 & 69.8 & 62.2 & 76.6 \\
    DINOv2 ViT-S/14 trainable                 & \cmark & \cmark & \cmark & \cmark & 86.0 & 74.4 & 66.0 & 68.9 & 75.7 & 86.4 & 72.1 & 63.2 & 56.0 & 72.5 \\
    DINOv2 ViT-S/14 nuSc rect., frozen        & \cmark & \cmark & \cmark & \cmark & 87.2 & 78.1 & 72.1 & 71.0 & 78.8 & 87.9 & 75.5 & 70.2 & 66.6 & 77.2 \\
    DINOv2 ViT-S/14 nuSc rect., trainable     & \cmark & \cmark & \cmark & \cmark & 88.5 & 79.7 & 74.1 & 67.9 & \textbf{79.7} & \textbf{88.3} & \textbf{79.8} & \textbf{73.1} & \textbf{71.3} & \textbf{79.9} \\
    DINOv2 ViT-S/14 nuSc sq., frozen          & \cmark & \cmark & \cmark & \cmark & 88.0 & \textbf{80.4} & \textbf{75.9} & 65.8 & \textbf{79.7} & 86.8 & 75.0 & 69.5 & 61.2 & 75.7 \\
    DINOv2 ViT-S/14 nuSc sq., trainable       & \cmark & \cmark & \cmark & \cmark & 87.7 & 77.1 & 68.4 & \textbf{71.5} & 78.0 & 83.7 & 73.0 & 64.2 & 66.7 & 73.8 \\
    \midrule
    MAE ViT-B/16 frozen                       & \cmark & \cmark & \cmark & \cmark & 89.5 & 80.9 & 75.1 & 65.5 & 80.2 & 87.6 & 76.9 & 70.5 & 69.5 & 78.0 \\
    MAE ViT-B/16 trainable                    & \cmark & \cmark & \cmark & \cmark & 86.3 & 78.6 & 69.3 & 67.8 & 77.6 & 86.3 & 76.5 & 68.5 & \textbf{74.9} & 77.9 \\
    MAE ViT-S/14 nuSc rect., frozen           & \cmark & \cmark & \cmark & \cmark & \textbf{90.3} & 80.2 & 74.2 & \textbf{73.7} & 81.4 & 85.3 & 77.3 & 72.9 & 69.9 & 77.9 \\
    MAE ViT-S/14 nuSc rect., trainable        & \cmark & \cmark & \cmark & \cmark & 89.9 & \textbf{82.4} & \textbf{77.4} & 72.6 & \textbf{82.4} & 89.2 & \textbf{79.9} & 70.2 & 68.9 & \textbf{79.3} \\
    MAE ViT-S/14 nuSc sq., frozen             & \cmark & \cmark & \cmark & \cmark & 88.9 & 78.6 & 73.7 & 68.3 & 79.4 & 84.5 & 75.2 & 71.5 & 67.3 & 76.3 \\
    MAE ViT-S/14 nuSc sq., trainable          & \cmark & \cmark & \cmark & \cmark & 87.4 & 77.8 & 74.2 & 70.6 & 79.2 & 85.8 & 77.7 & \textbf{74.4} & 64.1 & 77.6 \\
    \bottomrule
    \end{tabular}}%

    \vspace{0.6em}
    \footnotesize
    \emph{Single-city transfer rows for DINOv2 and MAE (Boston, Pittsburgh, Las~Vegas, Singapore) follow the same convention and are reported in the companion paper's supplementary material~\cite{naeinian-cross-city-2026}; the reduced set above is retained as the minimum provenance needed to justify the v1.1 re-evaluation used for all thesis claims.}
\end{sidewaystable}

% Full NAVSIM v1.1 cross-city PDMS matrix, canonical thesis table.
\begin{sidewaystable}[p]
    \centering
    \caption[NAVSIM v1.1 cross-city PDMS matrix (canonical).]{%
    Full closed-loop cross-city PDMS matrix on NAVSIM, devkit v1.1. This is the canonical closed-loop table used for every cross-city claim in the thesis. Each row is a single trained model; \emph{Train City} marks the cities used to fit that model. The two planner blocks (\emph{TransFuser}, \emph{Latent TransFuser}) report destination-city PDMS on the four NAVSIM test splits (LV, B, P, S) plus a \emph{Total} column averaging the four.
    How to read it: rows with all four training cities checked measure mixed-city interpolation; rows with a single training city measure zero-shot geographic transfer, in which case the three non-checked columns are the out-of-distribution cells. The gap between the in-distribution column and the OOD mean under a single planner isolates the transfer penalty; comparing the same row across the two planners separates backbone effects from head effects. All values are PDMS percentages; higher is better.}
    \label{tab:navsim_pdms_merged_v11}
    \scriptsize
    \setlength{\tabcolsep}{2pt}
    \renewcommand{\arraystretch}{0.95}
    \resizebox{\textwidth}{!}{%
    \begin{tabular}{|l|cccc|ccccc|ccccc|}
    \toprule
    \multirow{2}{*}{\textbf{Backbone}}
    & \multicolumn{4}{c|}{\textbf{Train City}}
    & \multicolumn{5}{c|}{\textbf{TransFuser (PDMS \%)}}
    & \multicolumn{5}{c|}{\textbf{Latent TransFuser (PDMS \%)}} \\
    \cmidrule(lr){2-5}\cmidrule(lr){6-10}\cmidrule(lr){11-15}
    & LV & B & P & S
    & LV & B & P & S & Total
    & LV & B & P & S & Total \\
    \midrule
    ResNet-34 supervised          & \cmark & \cmark & \cmark & \cmark & 86.6 & \textbf{78.4} & \textbf{73.8} & \textbf{64.0} & \textbf{77.9} & \textbf{89.0} & \textbf{82.5} & \textbf{80.3} & \textbf{66.5} & \textbf{81.7} \\
    ResNet-34 supervised          & --     & \cmark & --     & --     & 75.3 & 73.3 & 53.6 & 41.5 & 65.0 & 67.3 & 71.7 & 54.5 & 45.0 & 62.6 \\
    ResNet-34 supervised          & \cmark & --     & --     & --     & \textbf{89.2} & 67.0 & 55.1 & 45.3 & 68.6 & 86.9 & 62.1 & 53.2 & 42.3 & 65.5 \\
    ResNet-34 supervised          & --     & --     & \cmark & --     & 68.8 & 55.3 & 52.9 & 43.0 & 57.4 & 67.7 & 57.5 & 53.0 & 46.5 & 58.3 \\
    ResNet-34 supervised          & --     & --     & --     & \cmark & 56.6 & 45.4 & 49.6 & 42.1 & 49.5 & 57.4 & 42.6 & 45.1 & 45.5 & 48.5 \\
    \midrule
    I-JEPA ViT-H/14 frozen        & \cmark & \cmark & \cmark & \cmark & 85.3 & 75.7 & 72.4 & 67.3 & 76.9 & 83.7 & 75.2 & 70.9 & \textbf{68.3} & 76.1 \\
    I-JEPA ViT-H/14 trainable     & \cmark & \cmark & \cmark & \cmark & 86.6 & \textbf{77.1} & 70.6 & 64.8 & 77.0 & \textbf{88.7} & \textbf{77.3} & 75.1 & 67.1 & \textbf{79.1} \\
    I-JEPA ViT-S/14 nuSc rect., frozen        & \cmark & \cmark & \cmark & \cmark & 87.3 & 77.0 & \textbf{74.3} & \textbf{69.2} & \textbf{78.7} & 86.6 & 73.3 & 69.7 & 59.9 & 74.9 \\
    I-JEPA ViT-S/14 nuSc rect., trainable     & \cmark & \cmark & \cmark & \cmark & 89.0 & 76.8 & 73.3 & 67.0 & 78.6 & 86.2 & 73.9 & 74.7 & 59.5 & 75.9 \\
    I-JEPA ViT-S/14 nuSc sq., frozen          & \cmark & \cmark & \cmark & \cmark & 84.8 & 75.1 & 69.7 & 69.1 & 76.3 & 83.8 & 73.2 & 68.6 & 61.5 & 74.0 \\
    I-JEPA ViT-S/14 nuSc sq., trainable       & \cmark & \cmark & \cmark & \cmark & 85.8 & 75.7 & 73.0 & 65.2 & 76.9 & 87.3 & 76.4 & \textbf{75.3} & 67.7 & 78.5 \\
    I-JEPA ViT-H/14 frozen        & --     & \cmark & --     & --     & 72.5 & 69.9 & 58.2 & 50.1 & 65.3 & 69.4 & 67.5 & 62.2 & 42.2 & 63.1 \\
    I-JEPA ViT-H/14 trainable     & --     & \cmark & --     & --     & 72.0 & 68.9 & 60.1 & 42.0 & 63.9 & 70.7 & 68.2 & 60.7 & 48.7 & 64.5 \\
    I-JEPA ViT-H/14 frozen        & \cmark & --     & --     & --     & 86.9 & 61.2 & 53.7 & 51.4 & 66.7 & 85.9 & 57.5 & 51.9 & 47.8 & 64.3 \\
    I-JEPA ViT-H/14 trainable     & \cmark & --     & --     & --     & 87.7 & 60.1 & 53.9 & 51.3 & 66.7 & 82.7 & 60.3 & 54.8 & 44.3 & 64.2 \\
    I-JEPA ViT-H/14 frozen        & --     & --     & \cmark & --     & 71.4 & 61.3 & 64.5 & 49.5 & 63.5 & 71.8 & 63.5 & 69.0 & 48.6 & 65.0 \\
    I-JEPA ViT-H/14 trainable     & --     & --     & \cmark & --     & 69.1 & 59.6 & 65.3 & 39.0 & 60.7 & 63.7 & 60.1 & 59.4 & 35.4 & 57.3 \\
    I-JEPA ViT-H/14 frozen        & --     & --     & --     & \cmark & 36.0 & 35.9 & 45.1 & 35.6 & 37.7 & 58.3 & 44.6 & 43.0 & 55.6 & 50.6 \\
    I-JEPA ViT-H/14 trainable     & --     & --     & --     & \cmark & 56.7 & 45.3 & 48.0 & 50.3 & 50.4 & 62.7 & 43.9 & 44.4 & 52.8 & 51.7 \\
    \midrule
    DINOv2 ViT-S/14 frozen                    & \cmark & \cmark & \cmark & \cmark & 87.7 & 77.0 & 71.5 & 66.0 & 77.7 & 86.4 & 72.8 & 67.2 & 59.2 & 74.1 \\
    DINOv2 ViT-S/14 trainable                 & \cmark & \cmark & \cmark & \cmark & 85.4 & 73.5 & 67.6 & 68.9 & 75.6 & 82.7 & 68.9 & 64.0 & 54.2 & 70.2 \\
    DINOv2 ViT-S/14 nuSc rect., frozen        & \cmark & \cmark & \cmark & \cmark & 86.6 & 77.2 & 74.0 & 70.5 & 78.7 & 87.1 & 74.8 & 71.6 & \textbf{71.5} & 77.7 \\
    DINOv2 ViT-S/14 nuSc rect., trainable     & \cmark & \cmark & \cmark & \cmark & 87.9 & 79.0 & 76.2 & 67.4 & \textbf{79.6} & 87.3 & \textbf{78.5} & \textbf{75.7} & 66.7 & \textbf{79.0} \\
    DINOv2 ViT-S/14 nuSc sq., frozen          & \cmark & \cmark & \cmark & \cmark & 87.2 & \textbf{79.2} & \textbf{77.8} & 65.5 & 79.4 & 86.3 & 73.1 & 72.7 & 60.8 & 75.5 \\
    DINOv2 ViT-S/14 nuSc sq., trainable       & \cmark & \cmark & \cmark & \cmark & 86.3 & 75.9 & 69.0 & \textbf{71.2} & 77.3 & 86.7 & 73.6 & 70.6 & 71.2 & 77.0 \\
    \midrule
    MAE ViT-B/16 frozen                       & \cmark & \cmark & \cmark & \cmark & 89.0 & 80.8 & \textbf{78.7} & 65.1 & 80.7 & 81.7 & 70.7 & 62.8 & 69.9 & 72.7 \\
    MAE ViT-B/16 trainable                    & \cmark & \cmark & \cmark & \cmark & 85.4 & 78.0 & 73.9 & 67.4 & 78.0 & 88.7 & \textbf{81.9} & \textbf{80.0} & 61.6 & 80.6 \\
    MAE ViT-S/14 nuSc rect., frozen           & \cmark & \cmark & \cmark & \cmark & 89.2 & 79.1 & 75.1 & \textbf{73.1} & 80.7 & 87.4 & 79.2 & 79.8 & 64.6 & 79.8 \\
    MAE ViT-S/14 nuSc rect., trainable        & \cmark & \cmark & \cmark & \cmark & 89.3 & \textbf{81.2} & 78.4 & 72.1 & \textbf{81.9} & 86.2 & 79.7 & 77.8 & 72.3 & 80.3 \\
    MAE ViT-S/14 nuSc sq., frozen             & \cmark & \cmark & \cmark & \cmark & 87.7 & 77.3 & 75.2 & 67.8 & 78.9 & 86.4 & 75.0 & 70.3 & 67.9 & 76.7 \\
    MAE ViT-S/14 nuSc sq., trainable          & \cmark & \cmark & \cmark & \cmark & 86.9 & 77.3 & 76.3 & 70.3 & 79.2 & \textbf{89.3} & 78.5 & 78.0 & \textbf{74.1} & \textbf{81.4} \\
    MAE ViT-B/16 frozen                       & --     & \cmark & --     & --     & 70.6 & 67.0 & 56.8 & 43.4 & 62.4 & 71.8 & 73.4 & 60.0 & 48.7 & 66.3 \\
    MAE ViT-B/16 trainable                    & --     & \cmark & --     & --     & 70.6 & 67.3 & 58.0 & 41.9 & 62.5 & 69.4 & 73.1 & 61.3 & 53.7 & 66.5 \\
    MAE ViT-B/16 frozen                       & \cmark & --     & --     & --     & 88.3 & 66.2 & 55.2 & 52.9 & 69.3 & 86.6 & 66.0 & 56.1 & 52.2 & 68.8 \\
    MAE ViT-B/16 trainable                    & \cmark & --     & --     & --     & 87.4 & 59.8 & 53.2 & 51.6 & 66.4 & 87.8 & 63.3 & 56.0 & 51.0 & 68.1 \\
    MAE ViT-B/16 frozen                       & --     & --     & \cmark & --     & 69.3 & 59.9 & 61.1 & 48.9 & 61.5 & 71.2 & 54.8 & 50.4 & 52.6 & 59.1 \\
    MAE ViT-B/16 trainable                    & --     & --     & \cmark & --     & 73.0 & 65.6 & 68.8 & 40.8 & 64.8 & 64.4 & 63.9 & 61.0 & 35.7 & 59.1 \\
    MAE ViT-B/16 frozen                       & --     & --     & --     & \cmark & 58.0 & 45.3 & 51.2 & 56.7 & 52.5 & 54.6 & 44.0 & 46.6 & 49.7 & 49.0 \\
    MAE ViT-B/16 trainable                    & --     & --     & --     & \cmark & 51.0 & 46.0 & 52.2 & 36.0 & 47.3 & 55.6 & 43.8 & 48.0 & 42.5 & 48.4 \\
    \bottomrule
    \end{tabular}}
\end{sidewaystable}

\begin{table}[ht]
\centering
\caption[Boston-only v1.1 Latent TransFuser cross-city summary.]{%
Compact summary of the canonical Boston-trained Latent TransFuser rows from the v1.1 matrix in~\reftab{tab:navsim_pdms_merged_v11}. These four rows are used as the matched supervised-vs-SSL comparison throughout the thesis: each model is trained on Boston only and evaluated on every NAVSIM city. LV/B/P/S are destination-city PDMS percentages; the in-distribution column (bold) is Boston; \emph{OOD mean} averages the three remaining cities and \emph{Drop} is the in-distribution-minus-OOD-mean gap in percentage points. A smaller drop indicates a more geographically robust representation under a fixed planner. The MAE row has the best absolute OOD mean; the DINOv2 row has the smallest drop; the ResNet-34 row has the largest drop and is the supervised reference.}
\label{tab:navsim_pdms_boston_v11}
\vspace{0.35em}
\small
\renewcommand{\arraystretch}{1.15}
\setlength{\tabcolsep}{5pt}
\begin{tabular}{@{}>{\raggedright\arraybackslash}p{3.6cm} c c c c c c@{}}
\toprule
\textbf{Backbone} & \textbf{LV (\%)} & \textbf{B (\%)} & \textbf{P (\%)} & \textbf{S (\%)} & \textbf{OOD mean (\%)} & \textbf{Drop (pp)} \\
\midrule
\rowcolor{gray!12}
ResNet-34 supervised & 67.3 & \textbf{71.7} & 54.5 & 45.0 & 55.6 & \textbf{16.1} \\
MAE ViT-B/16 frozen  & 71.8 & \textbf{73.4} & 60.0 & 48.7 & \textbf{60.2} & 13.2 \\
I-JEPA ViT-H/14 frozen & 69.4 & \textbf{67.5} & 62.2 & 42.2 & 57.9 & 9.6 \\
DINOv2 ViT-S/14 frozen & 70.7 & \textbf{63.6} & 55.8 & 45.5 & 57.3 & \textbf{6.3} \\
\bottomrule
\end{tabular}
\end{table}

\FloatBarrier

\subsection{Best Training Ground (RQ2)}
\label{sec:exp-phase1-best-training-city}

Which source city gives the most transferable planner? \reffig{fig:iros-best-training-city} aggregates this across all backbones and heads. Las Vegas and Boston alternate as the top two; Pittsburgh is consistently mid-tier; Singapore is the worst single-city training source in every configuration, for the reasons discussed in the data card (\refapp{app:city-data-card}). The same pattern appears more compactly in the radar plot (\reffig{fig:iros-radar}), which shows per-city in-distribution and OOD PDMS for each of the strongest backbones. The radar reveals a useful operational criterion: the best single-city training source is the one whose in-distribution contour most nearly matches its OOD contour, i.e.\ whose model maintains similar competence under transfer rather than simply being strongest on its own city.

\begin{figure}[ht]
    \centering
    \includegraphics[width=0.75\textwidth]{iros/fig6_best_training_city.pdf}
    \caption{Average cross-city PDMS achieved by each of the four single-city training sources, grouped by backbone family. The ranking Las Vegas $\geq$ Boston $>$ Pittsburgh $>$ Singapore is stable across backbones; the absolute spread is larger for supervised ResNet-34 than for frozen SSL variants, consistent with the representation-centric account. Source: Fig.~6 of~\cite{naeinian-cross-city-2026}.}
    \label{fig:iros-best-training-city}
\end{figure}

\begin{figure}[ht]
    \centering
    \includegraphics[width=0.9\textwidth]{iros/fig12_radar_summary.pdf}
    \caption{Radar summary of per-city PDMS by training source (rows) and backbone family (panels). The filled area is the average of the four per-city scores; wider, rounder areas correspond to more uniform transfer. Frozen SSL radars are more nearly circular than supervised ResNet radars, which is a compact visual restatement of the order-of-magnitude robustness claim. Source: Fig.~12 of~\cite{naeinian-cross-city-2026}.}
    \label{fig:iros-radar}
\end{figure}

\subsection{Transfer Asymmetry Analysis}
\label{sec:exp-phase1-asymmetry}

Transfer asymmetry is not noise; it is one of the most important findings of the benchmark. The hardest direction is typically from a right-hand-traffic city into Singapore, not the reverse. That asymmetry is consistent with a representational explanation: a backbone trained only on right-hand-traffic scenes can overfit to layout priors that become actively misleading in Singapore, whereas a model trained in Singapore is forced to encode some structure that remains useful when moved back to right-hand-traffic cities. Dataset scale compounds the effect. Singapore is also the smallest source domain in NAVSIM, so a model trained there has fewer opportunities to learn a diverse but reusable representation~\cite{naeinian-cross-city-2026}.

\begin{figure}[ht]
    \centering
    \includegraphics[width=0.7\textwidth]{iros/fig8_asymmetry.pdf}
    \caption[Boston--Singapore transfer asymmetry (bar chart).]{Boston$\leftrightarrow$Singapore transfer asymmetry. Bars are OOD PDMS at the destination; the grouped comparison isolates the direction of transfer. Boston$\rightarrow$Singapore is the harsher direction for every backbone, consistent with the dataset-scale and traffic-convention asymmetry discussed in-text. The absolute asymmetry gap is smaller for frozen SSL than for supervised ResNet-34. Source: Fig.~8 of~\cite{naeinian-cross-city-2026}.}
    \label{fig:iros-asymmetry}
\end{figure}

\subsection{Per-Metric Breakdown}
\label{sec:exp-phase1-metrics}

The phase-1 preprint primarily reports aggregate PDMS rather than a full decomposition into the five PDMS components for every transfer pair. For that reason, the thesis uses total PDMS as the main closed-loop robustness indicator in this chapter. The important point is that the gap is visible at the level of overall driving competence before any sub-score decomposition is introduced: geographic shift degrades planning quality substantially, and the degradation is consistently smaller when the backbone is self-supervised.

\begin{figure}[ht]
    \centering
    \includegraphics[width=0.85\textwidth]{iros/fig14_allcity_training.pdf}
    \caption{All-city training comparison. Mixed-city NAVSIM training pools examples from every city, which narrows the gap between supervised and SSL backbones as the dataset becomes geographically diverse. The takeaway is not that supervised backbones are equivalent to SSL --- they remain behind on the transfer matrices --- but that any benchmark that only reports mixed-city PDMS is measuring a regime in which the representation question is effectively washed out. Source: Fig.~14 of~\cite{naeinian-cross-city-2026}.}
    \label{fig:iros-allcity}
\end{figure}

% ============================================================================
\section{SSL Backbones in Generative Planning}
\label{sec:exp-phase2}

The generative-planning study revisits in-distribution NAVSIM under a planner that is designed to attenuate backbone contributions: DiffusionDrive~\cite{jia-diffusiondrive-2025}. Chapter~\ref{chap:methodology}, Section~\ref{sec:meth-dd-why} develops the rationale in full; the short form is that DiffusionDrive denoises trajectories from K-means anchor initializations in two steps, conditions through grid-sampled BEV cross-attention, and uses time-only FiLM, and the composite effect is a head that can produce good trajectories even when the image features are weak. This makes DiffusionDrive a deliberately hostile substrate for the representation claim. The goal of this study is therefore not to top an in-distribution leaderboard; it is to measure whether any backbone-dependent structure survives DiffusionDrive's attenuation, because that is the harder precondition for the cross-city generative study in \refsec{sec:exp-phase3}. The ViT adapter is held identical to the one used in the cross-city benchmark (both pipelines import \texttt{create\_image\_encoder} from the same TransFuser backbone module), so the planner head is the only manipulated variable between the two studies.

\begin{table}[ht]
\centering
\caption{DiffusionDrive + SSL backbone results on NAVSIM (all cities).}
\label{tab:diffusiondrive-allcity}
\begin{tabular}{lllcc}
\toprule
\textbf{Exp} & \textbf{Backbone} & \textbf{Modality} & \textbf{Frozen} & \textbf{PDMS} \\
\midrule
A12   & ResNet34              & C+L & ---       & 87.9 \\
A12-o & MAE (rect, trainable) & C+L & No        & 86.9 \\
A12-u & MAE (sq, trainable)   & C+L & No        & 86.8 \\
A12-q & I-JEPA (sq, trainable)& C+L & No        & 86.2 \\
A10-j & I-JEPA nuSc (rect, frozen) & C+L & Yes  & 85.2 \\
A10-p & I-JEPA nuSc (sq, frozen)   & C+L & Yes  & 86.3 \\
A13   & ResNet34 (camera-only)& C   & ---       & 86.6 \\
\bottomrule
\end{tabular}
\end{table}

\subsection{SSL vs.\ Supervised in Generative Planning}
\label{sec:exp-phase2-ssl-vs-sup}

The $\sim$1pp gap between SSL (86.2--86.9) and supervised ResNet34 (87.9) on all-city training mirrors the pattern from the cross-city benchmark. SSL features are slightly weaker in-distribution in this architecture as well. That observation is important in the context of the broader thesis: if mixed-city PDMS were the only question, stronger planner design could easily dominate the narrative. Instead, the small in-distribution cost becomes interesting only because later studies ask whether SSL repays that cost under domain shift.

\subsection{Camera-Only: SSL Partially Replaces LiDAR}
\label{sec:exp-phase2-camera-only}

The camera-only baseline is already stronger than one might expect. ResNet34 camera-only (A13) reaches 86.6 PDMS, only 1.3 points below the corresponding camera+LiDAR configuration at 87.9. This narrows the remaining question: the issue is no longer whether camera-only planning is viable at all, but whether stronger visual pretraining can close the residual gap without reintroducing expensive sensors. The planned I-JEPA camera-only experiment is therefore not a side experiment; it is the most direct test of whether SSL features can substitute for part of the geometric information usually supplied by LiDAR.

\paragraph{Planned I-JEPA camera-only comparison (designed, pending).}
Starting from the A13 ResNet-34 camera-only baseline ($86.6$ PDMS, \reftab{tab:diffusiondrive-allcity}), the supervised encoder is replaced with the strongest frozen SSL variant: I-JEPA ViT-S/14 nuScenes square, frozen (A10-p, $86.3$ PDMS in a camera+LiDAR configuration). Configuration is in \reftab{tab:exp-camonly-config}; pre-registered reading bands are in \reftab{tab:exp-camonly-bands}.

\begin{table}[ht]
\centering
\caption{Configuration for the planned I-JEPA camera-only comparison. All fields are held fixed to the A10-p / A13 hyperparameters; only the backbone and the LiDAR modality differ from their respective cells of \reftab{tab:diffusiondrive-allcity}.}
\label{tab:exp-camonly-config}
\small
\setlength{\tabcolsep}{6pt}
\renewcommand{\arraystretch}{1.15}
\begin{tabular}{@{}>{\raggedright\arraybackslash}p{3.6cm} >{\raggedright\arraybackslash}p{8.4cm}@{}}
\toprule
\textbf{Setting} & \textbf{Value} \\
\midrule
Backbone        & I-JEPA ViT-S/14, nuScenes square pre-training, frozen \\
Modality        & Camera only (LiDAR removed) \\
Optimizer       & Adam, constant LR $1\times10^{-4}$ \\
Epochs          & $20$ \\
Anchor bank     & $256$ K-means trajectories on \texttt{navtrain} \\
Denoising steps & $2$ (truncated) at inference \\
Training split  & \texttt{navtrain}, all cities \\
Reference rows  & A10-p ($86.3$ PDMS, C+L), A13 ($86.6$ PDMS, camera-only) \\
\bottomrule
\end{tabular}
\end{table}

\begin{table}[ht]
\centering
\caption{Pre-registered reading of the planned camera-only PDMS outcome. Bands are fixed before the run; interpretation follows directly.}
\label{tab:exp-camonly-bands}
\small
\setlength{\tabcolsep}{6pt}
\renewcommand{\arraystretch}{1.2}
\begin{tabular}{@{}>{\raggedright\arraybackslash}p{2.5cm} >{\raggedright\arraybackslash}p{2.9cm} >{\raggedright\arraybackslash}p{6.6cm}@{}}
\toprule
\textbf{Band} & \textbf{PDMS range} & \textbf{Interpretation} \\
\midrule
Full compensation & $\geq 86.6$            & Representation fully replaces LiDAR \\
Partial           & $85.5$--$86.5$         & Modest gain over supervised camera-only at no LiDAR cost \\
Null              & $< 85.0$               & Frozen SSL does not recover the LiDAR contribution at this scale \\
\bottomrule
\end{tabular}
\end{table}

\subsection{Frozen vs.\ Fine-Tuned Backbone}
\label{sec:exp-phase2-frozen}

Within DiffusionDrive, the I-JEPA results show that preprocessing matters at least as much as fine-tuning. The rectangular frozen run reaches 85.2 PDMS, the square-crop frozen run reaches 86.3, and the trainable square-crop run reaches 86.2. The simplest interpretation is that matching the pretraining geometry of the ViT is more important than unfreezing the backbone, at least on the mixed-city split. This fits the broader thesis argument that frozen SSL weights act as a form of implicit regularization: they preserve the pretrained invariances while still allowing the adapter and planner to specialize downstream.

% ============================================================================
\section{Cross-City Generative Planning Protocol (Designed; Pending)}
\label{sec:exp-phase3}

The cross-city generative study extends the representation-centric cross-city benchmark into the generative-planning regime. The protocol is already defined and the training pipeline is ready to launch; the experimental runs themselves are scheduled outside the window of this thesis and will be reported in a follow-up publication. We document the protocol here with the same level of detail as the completed studies so that the design can be scrutinized (and replicated) independently of its execution.

\subsection{Rationale}
\label{sec:exp-phase3-rationale}

Two observations from the preceding studies motivate the cross-city generative study. First, Drive-JEPA~\cite{wang-drive-jepa} showed that mixed-city NAVSIM is no longer informative about a backbone's cross-city invariance --- a sufficiently strong generative planner closes the in-distribution gap regardless of backbone, which is precisely the effect our generative-planning results replicate for DiffusionDrive. Second, our earlier cross-city benchmark~\cite{naeinian-cross-city-2026} showed that geographic shift is the regime in which representation choice becomes load-bearing for the lightweight LAW and TransFuser heads.

The cross-city generative study asks whether that second observation survives the substitution of a generative planner for a regression one \emph{where the generative planner is hostile to the backbone by construction}. The reading criterion, set in advance and developed in \refsec{sec:meth-dd-success}, is directional: frozen SSL is required to produce a smaller OOD PDMS drop than supervised ResNet-34 in three or more of the four training-city rows for the representation-centric claim to be judged architecture-agnostic. Because DiffusionDrive's K-means anchor initialization, two-step denoising, and grid-sampled BEV conditioning attenuate the backbone's contribution, small absolute PDMS differences at the ratio level (in-distribution over OOD) are the load-bearing quantity; the thesis reports transfer-ratio inflation rather than raw PDMS gaps for this study.

\subsection{Experimental Matrix}
\label{sec:exp-phase3-matrix}

Two backbones $\times$ four training cities = eight training runs; each run is then evaluated on all four city test sets, for 32 closed-loop evaluations total. The backbones are chosen to pair a supervised baseline with the best generative-planning SSL configuration:
\begin{itemize}
    \item \textbf{ResNet-34 (C+L)} --- the all-city A12 configuration (87.9 PDMS, Table~\ref{tab:diffusiondrive-allcity}). Trainable, following the DiffusionDrive recipe.
    \item \textbf{I-JEPA ViT-S/14 nuScenes sq.\ frozen (C+L)} --- the A10-p configuration (86.3 PDMS), selected because (a) its square-crop geometry matches the ViT pretraining, (b) frozen SSL is the operational lever that matters for transfer per \refsec{sec:exp-why-ssl}, and (c) it is the strongest frozen DiffusionDrive configuration we measured in the generative-planning study.
\end{itemize}

Except for the training-city filter, every other hyperparameter (optimizer, learning rate, epoch count, batch size, anchor bank, denoising steps) is inherited from the generative-planning study's all-city configuration so that the cross-city generative study isolates the effect of geographic shift.

% Legacy protocol sketch retained for provenance. The completed matrix is in experiments.tex.
\begin{table}[ht]
    \centering
    \caption[Legacy cross-city generative study protocol sketch.]{Legacy protocol sketch for the initial two-backbone cross-city generative plan. The completed DiffusionDrive matrix reported in Chapter~\ref{chap:experiments} supersedes this table.}
    \label{tab:phase3-plan}
    \small
    \setlength{\tabcolsep}{3pt}
    \renewcommand{\arraystretch}{0.95}

    \begin{adjustbox}{max width=\textwidth}
    \begin{tabular}{|l|l|cccc|cccc|c|}
    \toprule
    \multirow{2}{*}{\textbf{Backbone}}
        & \multirow{2}{*}{\textbf{Run ID}}
        & \multicolumn{4}{c|}{\textbf{Train City}}
        & \multicolumn{5}{c|}{\textbf{DiffusionDrive PDMS (\%)}} \\
    \cmidrule(lr){3-6}\cmidrule(lr){7-11}
        & & LV & B & P & S & LV & B & P & S & Total \\
    \midrule
    ResNet-34 (C+L) & \texttt{DD-CC-Rn34-veg} & \cmark & -- & -- & --  & -- & -- & -- & -- & -- \\
    ResNet-34 (C+L) & \texttt{DD-CC-Rn34-bos} & -- & \cmark & -- & --  & -- & -- & -- & -- & -- \\
    ResNet-34 (C+L) & \texttt{DD-CC-Rn34-pit} & -- & -- & \cmark & --  & -- & -- & -- & -- & -- \\
    ResNet-34 (C+L) & \texttt{DD-CC-Rn34-sin} & -- & -- & -- & \cmark  & -- & -- & -- & -- & -- \\
    \midrule
    I-JEPA ViT-S/14 nuSc.\ rect.\ frozen (C+L) & \texttt{DD-CC-IJSr-veg} & \cmark & -- & -- & --  & -- & -- & -- & -- & -- \\
    I-JEPA ViT-S/14 nuSc.\ rect.\ frozen (C+L) & \texttt{DD-CC-IJSr-bos} & -- & \cmark & -- & --  & -- & -- & -- & -- & -- \\
    I-JEPA ViT-S/14 nuSc.\ rect.\ frozen (C+L) & \texttt{DD-CC-IJSr-pit} & -- & -- & \cmark & --  & -- & -- & -- & -- & -- \\
    I-JEPA ViT-S/14 nuSc.\ rect.\ frozen (C+L) & \texttt{DD-CC-IJSr-sin} & -- & -- & -- & \cmark  & -- & -- & -- & -- & -- \\
    \bottomrule
    \end{tabular}
    \end{adjustbox}
\end{table}


\subsection{Predicted Signatures and Falsifiable Tests}
\label{sec:exp-phase3-predictions}

The cross-city benchmark's v1.1 numbers ground the predictions below. The matched Boston-only Latent TransFuser rows from \reftab{tab:navsim_pdms_merged_v11} give an OOD PDMS drop of $16.1$ for trainable supervised ResNet-34 (F1-B1) against $9.6$ for frozen I-JEPA ViT-H/14 (F1-I1), a $\approx 2\times$ ratio that the mechanistic analysis (\refsec{sec:exp-implicit-reg}) links to feature-covariance effective rank within backbone families.

Pre-committed predictions are tabulated in \reftab{tab:exp-phase3-predictions}. The author will report cross-city generative numbers as they arrive and treat these bands as a falsifiable specification rather than a moving target.

\begin{table}[ht]
\centering
\caption{Pre-registered predictions for the cross-city generative study. Values derive from the matched LatTF rows in \reftab{tab:navsim_pdms_merged_v11} plus DiffusionDrive's $\sim 2$-point all-city uplift observed in the generative-planning study. ``OOD drop ratio'' is the ResNet-34 drop divided by the frozen I-JEPA drop.}
\label{tab:exp-phase3-predictions}
\small
\setlength{\tabcolsep}{5pt}
\renewcommand{\arraystretch}{1.2}
\begin{tabular}{@{}>{\raggedright\arraybackslash}p{4.4cm} >{\raggedright\arraybackslash}p{3.6cm} >{\raggedright\arraybackslash}p{4.6cm}@{}}
\toprule
\textbf{Prediction} & \textbf{Band} & \textbf{Falsification / reading rule} \\
\midrule
Single-city in-distribution PDMS, both backbones & $60$--$65$ & Values outside indicate the DD uplift or the LatTF baseline is misestimated \\
ResNet-34 OOD drop & $15$--$20$ PDMS & $<15$ weakens RQ5; $>20$ strengthens it \\
I-JEPA (frozen) OOD drop & $8$--$12$ PDMS & $<8$ strongly supports the representation claim \\
OOD drop ratio (Rn34/IJEPA) & $\approx 2\times$ & $<1.5\times$: weak evidence against; $>2.5\times$: strong confirmation \\
MAE-based DD (secondary) & Lowest OOD drop if added & From \reftab{tab:family-aggregate}: MAE frozen LatTF total $64.2$, drop $6.9$ \\
Singapore single-city row & Dominated by data scale & Separable from drop comparison by conditioning on LV / B / P \\
\bottomrule
\end{tabular}
\end{table}

\subsection{Link to the Mechanistic Analysis}
\label{sec:exp-phase3-mechanism}

The representation-level findings of \refsec{sec:exp-why-ssl} constrain the space of possible cross-city generative outcomes. If the mechanism is correctly identified as ``supervised-trainable backbones drift toward city-specific cues; frozen SSL backbones do not,'' then the cross-city generative study should produce a monotonic relationship between each run's probe-vs-OOD-drop slope and the planning head used to train it. Specifically, the predicted cross-city generative probe accuracies are: ResNet-34 DiffusionDrive $\approx$ 0.97--0.98 (same as the LAW/TransFuser numbers because the backbone is the same supervised ResNet), and I-JEPA-frozen DiffusionDrive $\approx$ 0.83 (same as F1-I1's 0.83 because the frozen SSL weights are identical regardless of the head trained on top). The planner does not touch the backbone; therefore the city-ID probe should report the same value.

If that probe prediction holds and the OOD-drop prediction above also holds, the representation claim is planner-invariant. If either fails, the mechanism as currently characterized is incomplete and the discussion in \refsec{sec:exp-why-synthesis} will need to admit a planner dependency.

\subsection{Scheduling Note}
\label{sec:exp-phase3-schedule}

The eight the cross-city generative study training jobs have been configured and queued in the same SLURM array driver that produced the mechanistic-analysis artifact set. The runs are gated on shared GPU availability on the NYU Greene cluster and are expected to complete in approximately one week of wall-clock time once they begin. Results will be reported in a follow-up manuscript; the present thesis closes on the cross-city, generative-planning, and mechanistic-analysis evidence.

\subsection{The ``1\% In-Distribution, 10\% OOD'' Rule}
\label{sec:exp-phase3-narrative}

Even before the cross-city generative study is populated, the project already supports a practical summary. On standard mixed-city training, SSL costs roughly a percentage point of PDMS relative to the strongest supervised all-city configuration (see the generative-planning gaps in \reftab{tab:diffusiondrive-allcity}). Under explicit city shift, that trade-off reverses: the representation becomes the dominant term, and frozen SSL backbones transfer far more stably (\reftab{tab:family-aggregate}). Accepting slightly lower mixed-city PDMS becomes a reasonable price to pay for substantially smaller geographic failure modes. the cross-city generative study will test whether this rule survives the planner swap; we record it here as the working summary of the thesis's practical claim.

% ============================================================================
\section{Why Does SSL Improve Driving Planning?}
\label{sec:exp-why-ssl}

The earlier phases establish the phenomenon. This section addresses the mechanism. The analyses below require no additional model training; they operate on cached features extracted from already trained checkpoints and ask a simple question: what is different about the feature space of SSL models that makes city transfer easier?

We probe 13 trained TransFuser / Latent TransFuser checkpoints (Tier-A) spanning supervised ResNet-34 and frozen / fully-trainable variants of I-JEPA, DINOv2, and MAE. For each (checkpoint, city) we extract global-average-pooled features from the deepest backbone stage on 556--1000 front-camera frames drawn from the NAVSIM test split of that city.\footnote{Feature dim is 512 for ResNet-34 and the TransFuser ViT-to-CNN adapter output for ViT backbones. The CLS token is dropped when present; spatial features are mean-pooled.} The five analyses below operate on these cached feature matrices; no additional training is required.

\subsection{Representation Similarity Across Cities (CKA)}
\label{sec:exp-cka}

Linear CKA is defined on two representations of the same samples, so we use the cross-dataset form that compares feature second-moment structures. For cities $A$ and $B$ with centered feature matrices $X_A$ and $X_B$, we report
\begin{equation}
\mathrm{CKA}(X_A, X_B)
=
\frac{\langle X_A^\top X_A,\; X_B^\top X_B \rangle_F}
     {\lVert X_A^\top X_A \rVert_F \,\lVert X_B^\top X_B \rVert_F},
\end{equation}
which is bounded in $[0, 1]$, invariant to orthogonal rotations of the feature basis, and agnostic to different sample counts across cities.

Mean off-diagonal CKA across the four cities is uniformly high --- family-level averages range from 0.74 (ResNet-34) through 0.77 (MAE frozen) and 0.79--0.83 (I-JEPA, DINOv2 frozen) to 0.86 (DINOv2 trainable TF) --- and does not cleanly separate supervised from SSL backbones. Both families encode similar covariance structure across cities, so covariance alignment alone is not the discriminating mechanism. Per-run heatmaps (Fig.~\ref{fig:cka-grid}) show city-pair asymmetries consistent with the raw scenario distributions: Boston $\leftrightarrow$ Pittsburgh are closest, Singapore is the outlier for every backbone.

\begin{figure}[t]
    \centering
    \includegraphics[width=0.45\textwidth]{phase4/cka_F1-B1.pdf}
    \includegraphics[width=0.45\textwidth]{phase4/cka_P1-I1.pdf}
    \caption{Cross-city linear-CKA heatmaps for Boston-only supervised ResNet-34 LatTF (F1-B1, left) vs.\ Boston-only I-JEPA ViT-H/14 frozen TF (P1-I1, right). All 188 per-run heatmaps are in the supplementary material (\texttt{fig/phase4/heatmaps/}).}
    \label{fig:cka-grid}
\end{figure}

\subsection{Linear Probing for City Classification}
\label{sec:exp-linear-probe}

We fit a multinomial logistic regression on standardized features to predict the four-way city label, with an 80/20 stratified split. A random-Gaussian feature negative control yields 25.4\% accuracy (near chance), confirming the probe is well calibrated.

Family-level probe accuracies (Table~\ref{tab:family-aggregate}) separate every backbone family from supervised ResNet-34 by a wide margin. Trainable supervised ResNet-34 encodes city identity almost losslessly (0.97--0.98), MAE frozen sits in the middle (0.89--0.90), I-JEPA frozen next (0.83), and DINOv2 frozen encodes city identity least (0.74--0.80). The single outlier toward feature collapse is A6-b (I-JEPA ViT-H/14 fully-trainable all-city TF), which probes at 0.38 --- an order of magnitude below its frozen counterpart A6 at 0.94. A second collapse at A7-b (DINOv2 all-city fully-trainable TF, probe 0.39) shows that the pathology is not unique to I-JEPA; Sec.~\ref{sec:exp-implicit-reg} quantifies the underlying feature-rank collapse.

The paired single-city comparison on Boston makes the representation difference concrete: running the same linear probe on features from four Boston-only models produces 0.97 for ResNet-34 (F1-B1), 0.89 for MAE (F1-M1), 0.83 for I-JEPA (F1-I1), and 0.81 for DINOv2 (F1-D1). The OOD PDMS drops from Table~\ref{tab:navsim_pdms_merged} on the same four runs are 16.1, 13.2, 9.6, and 6.3, respectively --- a monotonic ordering: the less city-specific the frozen features are, the more gracefully the trained planner transfers.

\begin{table}[h]
\centering
\small
\begin{tabular}{llccc}
\toprule
Run & Backbone & Probe acc. & Train PDMS & OOD drop \\
\midrule
F1-B1 & ResNet-34 LatTF (trainable)       & \textbf{0.968} & 71.7 & \textbf{16.1} \\
F1-M1 & MAE ViT-B/16 LatTF frozen         & 0.889 & 73.4 & 13.2 \\
F1-I1 & I-JEPA ViT-H/14 LatTF frozen      & 0.833 & 67.5 & 9.6 \\
F1-D1 & DINOv2 ViT-S/14 LatTF frozen      & \textbf{0.812} & 63.6 & \textbf{6.3} \\
\bottomrule
\end{tabular}
\caption{Boston-only Latent-TransFuser runs: linear-probe city-ID accuracy orders the four backbones the same way as their measured OOD PDMS drop. The trainable supervised ResNet-34 drifts most toward city-specific cues and drops most under geographic transfer; frozen DINOv2 is least city-discriminative and transfers best (Table~\ref{tab:navsim_pdms_merged}).}
\label{tab:probe-boston}
\end{table}

Across the full 188-run roster the story is more nuanced. Pooling every (backbone, variant, city) combination, Spearman $\rho$ between probe accuracy and OOD PDMS drop is only $-0.13$; stratified by training city it ranges from $-0.46$ (Pittsburgh-trained) to $+0.17$ (Boston-trained). The ``probe predicts drop'' law holds \emph{within} matched comparisons (same train-city, same backbone family, frozen vs.\ trainable; or same train-city across backbones) but washes out once training city, backbone scale, and pretraining data vary together. Fig.~\ref{fig:probe-vs-ood-stratified} breaks the scatter out per training city; Fig.~\ref{fig:probe-vs-ood} shows the pooled view for reference.

\begin{figure}[t]
    \centering
    \includegraphics[width=\textwidth]{phase4/scatter_probe_vs_ood_by_training_city.pdf}
    \caption{Linear-probe city-ID accuracy vs.\ train-city minus test-city PDMS, stratified by training city. Within each panel the expected ordering (higher probe $\rightarrow$ larger OOD drop) is visible for the Pittsburgh, Boston, and Las Vegas panels and flat for Singapore; pooling across training cities washes the signal out.}
    \label{fig:probe-vs-ood-stratified}
\end{figure}

\begin{figure}[t]
    \centering
    \includegraphics[width=0.65\textwidth]{phase4/scatter_probe_vs_ood.pdf}
    \caption{Pooled linear-probe accuracy vs.\ OOD PDMS drop across all 188 analyzed Tier-A runs. The global correlation ($\rho = -0.13$) is weak because training city dominates; the matched-comparison signal surfaces in Fig.~\ref{fig:probe-vs-ood-stratified} and Table~\ref{tab:probe-boston}.}
    \label{fig:probe-vs-ood}
\end{figure}

\subsection{Feature Visualization via Attention Maps}
\label{sec:exp-attention}

Last-block CLS-to-patch attention visualization (\texttt{attention\_viz.py}) requires a curated set of semantically matched scenes (one per city) to be visually informative. Scene curation is left as future work; the quantitative probe, CKA, and effective-rank results already carry the mechanistic argument.

\subsection{Feature Distribution Shift (MMD)}
\label{sec:exp-mmd}

We compute biased RBF-kernel MMD\textsuperscript{2} between city-pair feature distributions with a median-heuristic bandwidth $\sigma^2$ estimated from a shared 500-sample pool, clamping all squared distances at 0 to avoid numerical blow-up on near-duplicate samples (observed on a handful of trainable single-city DINOv2 runs pre-clamp). Absolute MMD\textsuperscript{2} values are not directly comparable across backbones because $\sigma^2$ differs; within-regime comparisons remain informative. Fig.~\ref{fig:mmd-vs-ood} is the pooled scatter; the Boston-only cluster reproduces the single-city ordering (Table~\ref{tab:probe-boston}).

Family-level mean MMD\textsuperscript{2} across all 188 runs: MAE 0.09--0.11, ResNet-34 0.10, I-JEPA 0.07--0.09, DINOv2 0.05--0.07. Pooled $\rho$(log MMD\textsuperscript{2}, OOD drop) is $-0.22$, weakly negative --- higher feature-space variance correlates with \emph{smaller} OOD drops. The reason is that MAE and ResNet-34, which produce the highest-variance features, also generalize best in their respective frozen and supervised regimes; collapsed-feature DINOv2 and I-JEPA runs give tight distributions but poor transfer. MMD is therefore a sanity check rather than a predictor.

\begin{figure}[t]
    \centering
    \includegraphics[width=0.65\textwidth]{phase4/scatter_mmd_vs_ood.pdf}
    \caption{$\log_{10}$ mean cross-city MMD\textsuperscript{2} vs.\ OOD PDMS drop over the full 188-run roster. Single-city points dominate the upper half; all-city points (drop $\approx 0$) sit at the bottom. Within the single-city cluster the MAE and ResNet-34 families occupy the higher-MMD region, DINOv2 the lower-MMD region.}
    \label{fig:mmd-vs-ood}
\end{figure}

\subsection{Feature-Covariance Rank (Implicit Regularization)}
\label{sec:exp-implicit-reg}

The original plan was to measure the effective rank of learned planning-head weights. With frozen SSL backbones the planning head is the only drift vector, which would make it a clean test; however, for supervised ResNet-34 the backbone itself is trainable, so planning-head rank alone cannot be compared fairly. We instead report the effective rank of the \emph{feature covariance} for each (run, city) pair: per-city features are centered, $\sigma_i$ denotes the $i$-th singular value of the centered matrix, and we report $\exp\!\bigl(-\sum_i (\tilde\sigma_i^2) \log (\tilde\sigma_i^2)\bigr)$ with $\tilde\sigma_i$ the $\ell_2$-normalized singular values~\cite{roy-effrank-2007}. This probes representational structure rather than downstream-head capacity.

Family-level effective ranks (Table~\ref{tab:family-aggregate}) separate the SSL objectives cleanly. MAE frozen TF retains the highest effective rank at 80.3, larger than trainable ResNet-34 (69.7) or I-JEPA frozen (66.2--69.0), with DINOv2 frozen a notch below (58.6--61.5). Fine-tuning reduces rank universally --- MAE trainable collapses to 47--49, I-JEPA trainable to 47--53, DINOv2 trainable stays at 52--65. The three most extreme collapses across the whole roster (effective rank $<5$) all involve trainable or single-city runs on large pretrained ViTs: A7-b (DINOv2 ViT-S trainable all-city TF, rank 1.2), A6-b (I-JEPA ViT-H trainable all-city TF, rank 2.9), and P1-D2 (DINOv2 frozen LV-only TF, rank 3.4). Feature-rank collapse is therefore not a quirk of one I-JEPA run; it is a systemic risk of fine-tuning high-capacity SSL backbones on modest driving datasets.

Across the 188-run pool, feature-covariance rank does \emph{not} monotonically predict OOD drop ($\rho = +0.08$). The informative read is within-family: MAE frozen LatTF has both the highest rank (72.6) and the smallest OOD drop (6.9); MAE trainable LatTF has half the rank (49.1) and a larger drop (9.2); I-JEPA follows the same pattern (frozen rank 69, drop 11.6; trainable rank 47, drop 9.5). The frozen constraint preserves the pretrained rank structure and, for MAE especially, keeps OOD transfer tight.

\begin{table}[h]
\centering
\footnotesize
\begin{tabular}{lllrrrrrr}
\toprule
Backbone & Status & Mode & $N$ & Probe & CKA & MMD$^2$ & Eff.\ rank & OOD drop \\
\midrule
ResNet-34 & supervised & TF    & 2  & 0.98 & 0.75 & 0.095 & 69.7 & $-4.2$\textsuperscript{$\dagger$} \\
ResNet-34 & supervised & LatTF & 5  & 0.98 & 0.75 & 0.097 & 63.3 & 8.7 \\
I-JEPA    & frozen     & TF    & 15 & 0.83 & 0.79 & 0.084 & 66.2 & 11.4 \\
I-JEPA    & frozen     & LatTF & 15 & 0.83 & 0.77 & 0.081 & 69.0 & 11.6 \\
I-JEPA    & trainable  & TF    & 15 & 0.75 & 0.82 & 0.067 & 52.8 & 12.5 \\
I-JEPA    & trainable  & LatTF & 15 & 0.83 & 0.83 & 0.089 & 46.5 & 9.5 \\
DINOv2    & frozen     & TF    & 15 & 0.74 & 0.85 & 0.059 & 61.5 & 12.4 \\
DINOv2    & frozen     & LatTF & 15 & 0.80 & 0.83 & 0.069 & 58.6 & 10.0 \\
DINOv2    & trainable  & TF    & 15 & 0.69 & 0.86 & 0.047 & 51.6 & 12.8 \\
DINOv2    & trainable  & LatTF & 15 & 0.77 & 0.82 & 0.062 & 65.3 & 10.0 \\
MAE       & frozen     & TF    & 15 & 0.89 & 0.75 & 0.094 & \textbf{80.3} & 11.2 \\
MAE       & frozen     & LatTF & 15 & 0.90 & 0.77 & 0.092 & 72.6 & \textbf{6.9} \\
MAE       & trainable  & TF    & 15 & 0.85 & 0.76 & 0.107 & 47.5 & 11.7 \\
MAE       & trainable  & LatTF & 15 & 0.89 & 0.78 & 0.097 & 49.1 & 9.2 \\
\bottomrule
\end{tabular}
\caption{Family-level aggregates averaged across every analyzed run in each cell. Probe = linear-probe city-ID accuracy; CKA = mean off-diagonal cross-city linear CKA; MMD$^2$ = RBF median-bandwidth biased MMD$^2$ averaged across city pairs; Eff.\ rank = Roy--Vetterli effective rank of the per-city feature covariance averaged across the four cities; OOD drop = train-city minus test-city PDMS (from Table~\ref{tab:navsim_pdms_merged}). \textsuperscript{$\dagger$}Only two ResNet-34 TF checkpoints were analyzed (P1-B1-b and P1-B2-b pointers are missing on disk); the negative drop is dominated by Pittsburgh- and Singapore-trained runs where in-distribution PDMS is pathologically low. Full per-run metrics are committed as \texttt{fig/phase4/summary.tsv} and \texttt{all\_metrics.json}.}
\label{tab:family-aggregate}
\end{table}

\begin{figure}[t]
    \centering
    \includegraphics[width=0.45\textwidth]{phase4/spectra/spectrum_A6.pdf}
    \includegraphics[width=0.45\textwidth]{phase4/spectra/spectrum_A6-b.pdf}
    \caption{Per-city normalized singular-value spectrum for I-JEPA ViT-H/14 frozen (A6, left) vs.\ fully-trainable (A6-b, right), both trained on all four cities. Full fine-tune of the ViT-H encoder compresses the feature spectrum by more than two orders of magnitude; the resulting near-rank-2 representation drives the corresponding probe accuracy down to 0.38.}
    \label{fig:spectrum-collapse}
\end{figure}

\subsection{Synthesis: A Mechanistic Account}
\label{sec:exp-why-synthesis}

The five analyses converge on three related findings.

First, trainable supervised backbones drift toward city-discriminative cues. A linear probe on trained ResNet-34 features separates the four cities at 0.97--0.98 accuracy across every configuration we trained, versus 0.25 for random features. Frozen SSL encoders cannot drift under driving gradients; their probe accuracies are an order of magnitude further from saturation (DINOv2 0.74--0.80, I-JEPA 0.83, MAE 0.89--0.90). In matched single-city comparisons --- one training city, one planning head, four different backbones --- the ordering of probe accuracies matches the ordering of OOD PDMS drops measured in Table~\ref{tab:navsim_pdms_merged} (Table~\ref{tab:probe-boston}). This is the representation-mechanism claim of the thesis, restricted to the cell in which it is identifiable.

Second, across the full 188-run roster the simple probe-vs-drop correlation is weak ($\rho = -0.13$, Fig.~\ref{fig:probe-vs-ood}) because training-city confounds dominate and because the MAE family --- the OOD winner by several points --- has moderately high probe accuracy without paying a generalization cost. MAE frozen LatTF (OOD drop 6.9, total PDMS 64.2) is the best family in the roster despite probing higher than DINOv2 and I-JEPA. The pretraining objective therefore matters beyond its effect on probe accuracy: the masked-autoencoder objective produces features that are both city-discriminable and city-robust, because it optimizes for reconstruction of scene content rather than for an embedding distance that can be minimized by discarding structure.

Third, the frozen constraint is itself load-bearing. Fine-tuning compresses feature-covariance rank for every SSL family (MAE 72$\rightarrow$47, I-JEPA 69$\rightarrow$47, DINOv2 60$\rightarrow$52--65), and in the worst cases produces an outright collapse to effective rank $<5$ (A6-b I-JEPA, A7-b DINOv2, P1-D2 DINOv2). The collapsed runs lose their probe-accuracy and MMD signals simultaneously --- with only two or three nontrivial feature directions left, the linear probe cannot separate cities and the RBF kernel cannot resolve distributions. Preserving the pretrained rank structure is, empirically, what keeps the SSL representation useful for transfer; the frozen SSL backbones that retain it (MAE ViT-B/16, I-JEPA ViT-H/14, DINOv2 ViT-S/14) are the ones that generalize.

Cross-city robustness in this study is therefore best summarized as three claims:
(i) supervised backbones learn city identity almost for free, and the fraction they learn predicts their transfer deficit within matched comparisons;
(ii) SSL objectives differ materially, with MAE leading the OOD rankings and DINOv2 most prone to collapse under fine-tune;
(iii) the operational lever with the highest leverage is not the SSL objective alone but whether the encoder is frozen at inference time.

% ============================================================================
\section{Ablation Studies}
\label{sec:exp-ablations}

\subsection{Input Resolution: Rectangular vs.\ Square Crop}
One ablation that is already informative from the generative-planning study is input geometry. For I-JEPA in DiffusionDrive, the square-crop frozen configuration (86.3 PDMS) outperforms the rectangular frozen configuration (85.2 PDMS). This suggests that matching the input geometry of ViT pretraining can matter more than preserving the raw driving aspect ratio, at least when the downstream adapter is lightweight.

\subsection{Frozen vs.\ Fine-Tuned Backbone (Extended)}
The frozen-versus-fine-tuned question remains central across studies. Available evidence from the generative-planning study and the mechanistic analysis suggests that fine-tuning can slightly improve in-distribution fit, but also risks eroding the pretrained invariances that make SSL useful under transfer. The feature-rank collapse observed in fully trainable I-JEPA provides a concrete mechanistic warning against unconstrained fine-tuning on modest driving datasets.

\subsection{Pre-Training Data: ImageNet vs.\ nuScenes}
The ImageNet-versus-driving-data comparison is addressed most directly in the cross-city benchmark, where domain-specific nuScenes pretraining consistently yields the most stable cross-city results. The lesson is not that generic pretraining is useless, but that domain alignment matters much more under geographic transfer than it does on mixed-city validation.

\subsection{ViT Scale: ViT-S vs.\ ViT-H}
Backbone scale is a practical trade-off rather than a purely academic one. The larger ViT-H variants are useful for diagnosing transfer in the cross-city benchmark, while smaller ViT-S variants are more realistic inside DiffusionDrive where throughput matters. The broader pattern in this thesis is that objective and freezing strategy matter more reliably for robustness than raw backbone size alone.

% ============================================================================
\section{EPDMS Evaluation}
\label{sec:exp-epdms}

Where available, EPDMS is used to place the all-city DiffusionDrive runs in the context of recent generative planners. The most important role of this section is comparative rather than leaderboard-oriented: it shows how far the thesis models remain from the strongest dedicated in-distribution systems such as Drive-JEPA, and therefore clarifies why the main claim of the thesis is about robustness rather than pure aggregate benchmark dominance.

% ============================================================================
\section{RL Post-Training Transfer (Exploratory)}
\label{sec:exp-rl}

RL post-training remains exploratory in this thesis. The intended test is to improve the best all-city JEPA-Diffusion model with reward-guided or preference-guided refinement and then re-evaluate it on held-out cities. The result would be interesting in either direction:
\begin{itemize}
    \item If RL gains transfer, then SSL + diffusion + RL becomes a plausible recipe for robust high-performance planning.
    \item If RL gains do not transfer, then reward optimization is likely overfitting to city-specific structure, which would itself be an important negative result.
\end{itemize}

% ============================================================================
\section{Qualitative Analysis}
\label{sec:exp-qualitative}

Qualitative analysis is most valuable in this thesis when it is used to illustrate failure modes that are already visible in the quantitative results. The most informative examples are:
\begin{itemize}
    \item success cases where SSL backbones preserve lane-following or collision avoidance under cross-city transfer while supervised baselines fail,
    \item failure cases where all backbones struggle, such as rare intersection geometries or highly ambiguous traffic interactions,
    \item and, for diffusion planners, examples of trajectory diversity and proposal selection under different backbone choices.
\end{itemize}

% ============================================================================
\section{Discussion}
\label{sec:exp-discussion}

Taken together, the studies support a single thesis argument. The exploratory pilots showed that SSL could already make lightweight planners competitive, which justified treating representation quality as a first-class variable rather than a background choice. The cross-city benchmark then made the key point visible: once evaluation is city-disjoint, backbone choice dominates robustness. The generative-planning study showed that this representation-centric story survives a major planner change in-distribution, and the mechanistic analysis explained why: frozen SSL features are less city-discriminative and occupy a more stable feature subspace than supervised or fully fine-tuned alternatives.

The practical recommendation is therefore asymmetric. If the target deployment environment is well matched to the training mix, the strongest supervised in-distribution planner may be adequate. If deployment includes genuinely new cities, the better trade-off is usually to accept a small in-distribution cost in exchange for much smaller geographic failure modes. The main limitations of the current evidence are the restriction to four cities, the heavy reliance on open-loop proxy evaluation, and the fact that the full DiffusionDrive cross-city matrix is still a designed extension rather than a complete empirical result.
